\chapter{Complexity — PTM Equiv}
%%% generated by the blueprint pipeline from TCSlib.Complexity.Randomized.PTMEquiv
\section{Overview}

Declarations of the Lean module \texttt{TCSlib.Complexity.Randomized.PTMEquiv}.

\section{Declarations}

\begin{theorem}[Length selected by a Boolean mask]
\lean{Randomized.length_maskSelect}
\label{Randomized.length_maskSelect}
\leanok
\uses{Complexity.maskSelect}
\difficulty{4}
Let $m$ and $w$ be Boolean lists of equal length. Select from $w$ the entries whose
corresponding entries in $m$ are true. The resulting list has length equal to the number
of true entries in $m$.
\end{theorem}

\begin{theorem}[Probability under Boolean mask selection]
\lean{Randomized.randProb_maskSelect}
\label{Randomized.randProb_maskSelect}
\leanok
\uses{Complexity.maskSelect, Randomized.randProb, Randomized.randProb_congr, Randomized.randProb_succ}
\difficulty{5}
Let $M$ be a finite Boolean mask and let $A$ be any predicate on finite Boolean words.
Apply the mask to a uniformly random Boolean word of length equal to the mask length,
retaining precisely the bits whose mask entries are true. The probability that the
resulting word satisfies $A$ equals the probability that a uniformly random Boolean word
of length equal to the number of true entries in the mask satisfies $A$.
\end{theorem}

\begin{theorem}[Probability under masked selection]
\lean{Randomized.randProb_maskSelect_drop_take}
\label{Randomized.randProb_maskSelect_drop_take}
\leanok
\uses{Complexity.maskSelect, Randomized.randProb, Randomized.randProb_drop_take, Randomized.randProb_maskSelect}
\difficulty{1}
Let $s,m$ be nonnegative integers, let $w$ be a uniformly random Boolean string of
length $m$, and let $A$ be a property of Boolean strings. For a Boolean mask whose
length satisfies $s+|\mathrm{mask}|\le m$, the probability that selecting from the
length-$|\mathrm{mask}|$ segment of $w$ immediately after its first $s$ bits, keeping
exactly the positions marked true, produces a string satisfying $A$ equals the
probability that a uniformly random string of length equal to the number of true entries
of the mask satisfies $A$.
\end{theorem}

\begin{theorem}[Boolean indicator of membership]
\lean{Randomized.indicator_eq_true_iff}
\label{Randomized.indicator_eq_true_iff}
\leanok
\uses{Turing.MultiTapeTM.indicator}
\difficulty{1}
For every language $V$ of bit strings and every bit string $y$, the Boolean indicator of
$V$ at $y$ is true if and only if $y\in V$.
\end{theorem}

\begin{theorem}[Boolean indicator of nonmembership]
\lean{Randomized.indicator_eq_false_iff}
\label{Randomized.indicator_eq_false_iff}
\leanok
\uses{Turing.MultiTapeTM.indicator}
\difficulty{1}
For every language $V$ of bit strings and every bit string $y$, the Boolean indicator of
$V$ at $y$ is false if and only if $y\notin V$.
\end{theorem}

\begin{theorem}[Efficiency from a polynomial-time test]
\lean{Randomized.eff_boolVerifier_of_test}
\label{Randomized.eff_boolVerifier_of_test}
\leanok
\uses{Complexity.PolyTimeComputable, Complexity.compl_mem_P, Complexity.mem_P_of_test, Randomized.boolVerifier, Randomized.polyTimeModel, Turing.pairEncode}
\difficulty{3}
Let $M$ be a Boolean-valued function of two finite Boolean strings, and let $t$ be a
Boolean-valued function on finite Boolean strings. Suppose the function that maps each
string $z$ to the one-symbol output containing $t(z)$ is computable in polynomial time,
and suppose $t(\langle x,r\rangle)=M(x,r)$ for every pair of strings, where $\langle
x,r\rangle$ is formed by doubling each bit of $x$, appending the separator $01$, and
then appending $r$. Then the verifier that accepts exactly when $M(x,r)$ is true is
efficient in the polynomial-time verifier model.
\end{theorem}

\begin{theorem}[Efficiency of an appended-input verifier]
\lean{Randomized.eff_append_mem}
\label{Randomized.eff_append_mem}
\leanok
\uses{Complexity.P, Complexity.PolyTimeComputable, Complexity.PolyTimeComputable.append, Complexity.compl_mem_P, Complexity.pairFstD, Complexity.pairFstD_pairEncode, Complexity.pairSndD, Complexity.pairSndD_pairEncode, Complexity.polyTimeComputable_pairFstD, Complexity.polyTimeComputable_pairSndD, Complexity.preimage_mem_P, Randomized.boolVerifier, Randomized.indicator_eq_false_iff, Randomized.indicator_eq_true_iff, Randomized.polyTimeModel, Turing.MultiTapeTM.indicator, Turing.pairEncode}
\difficulty{4}
Let $V$ be a language of finite Boolean strings that is decidable in polynomial time.
The verifier that, on input strings $x$ and $r$, accepts exactly when the concatenation
$x r$ belongs to $V$ is efficient in the polynomial-time verifier model.
\end{theorem}

\begin{theorem}[Membership determined by the truncated run]
\lean{Randomized.append_mem_iff_outputWith_take}
\label{Randomized.append_mem_iff_outputWith_take}
\leanok
\uses{Turing.FinNDTM, Turing.FinNDTM.outputWith, Turing.FinNDTM.runWith_eq_runWith_take, Turing.NDTM.HaltsWithin, Turing.NDTM.initCfg, Turing.NDTM.runWith}
\difficulty{4}
Let $N$ be a finite nondeterministic Turing machine, $V$ a language of bit strings, and
$x$ an input. Suppose that for every string $u$ of length $Q$, $x$ followed by $u$
belongs to $V$ exactly when the run of $N$ on $x$ under choices $u$ halts and outputs
the single-bit string true. If every run on $x$ under a choice string of length $T$
halts and $T\le Q$, then for every choice string $r$ of length $Q$, $x$ followed by $r$
belongs to $V$ exactly when the run under the first $T$ bits of $r$ outputs true.
\end{theorem}

\begin{theorem}[Choice verification and output probability]
\lean{Randomized.randProb_choiceVerifier}
\label{Randomized.randProb_choiceVerifier}
\leanok
\uses{Randomized.append_mem_iff_outputWith_take, Randomized.indicator_eq_true_iff, Randomized.randProb, Randomized.randProb_congr, Randomized.randProb_take, Turing.FinNDTM, Turing.FinNDTM.outputProb, Turing.FinNDTM.outputWith, Turing.MultiTapeTM.indicator, Turing.NDTM.HaltsWithin, Turing.NDTM.initCfg, Turing.NDTM.runWith}
\difficulty{3}
Let $N$ be a finite nondeterministic Turing machine over the Boolean alphabet, let $V$
be a language of Boolean words, and let $x$ be a Boolean word. Suppose that for every
word $u$ of length exactly $Q$, the concatenation $xu$ belongs to $V$ if and only if the
run of $N$ on input $x$ under choices $u$ halts and outputs the one-symbol word
$[\mathrm{true}]$. Suppose also that every run of $N$ on $x$ under a choice word of
length exactly $T$ halts, and that $T\le Q$. Then the probability that $xu$ belongs to
$V$ for a uniformly random Boolean word $u$ of length $Q$ equals the probability that
$N$ outputs $[\mathrm{true}]$ on input $x$ under a uniformly random choice word of
length $T$.
\end{theorem}

\begin{definition}[Bounded-error decision with a success threshold]
\lean{Turing.FinNDTM.BPDecidesInTimeWith}
\statementsource{Vad12}{pdf-bf2ce5422088-p023-b001}
\statementsource{arora-barak-ch07-randomized-computation}{pdf-26d95080dbde-p002-b003}
\label{Turing.FinNDTM.BPDecidesInTimeWith}
\leanok
\uses{Turing.FinNDTM, Turing.FinNDTM.outputProb, Turing.MultiTapeTM.indicator, Turing.NDTM.HaltsWithin}
A finite-state nondeterministic Turing machine $N$ decides a language $L$ within time
bound $T$ with success threshold $p$ if, for every Boolean input $x$, every computation
halts within $T(|x|)$ steps and the probability that its output equals the indicator of
membership of $x$ in $L$ is at least $p(|x|)$.
\end{definition}

\begin{theorem}[Two formulations of bounded-error decision]
\lean{Turing.FinNDTM.bpDecidesInTime_iff_with}
\label{Turing.FinNDTM.bpDecidesInTime_iff_with}
\leanok
\uses{Turing.FinNDTM, Turing.FinNDTM.BPDecidesInTime, Turing.FinNDTM.BPDecidesInTimeWith}
\difficulty{1}
Let $N$ be a finite nondeterministic Turing machine with Boolean input and output, let
$L\subseteq\{0,1\}^*$ be a language, and let $T:\mathbb{N}\to\mathbb{N}$. The machine
halts within time $T$ and outputs the indicator of $L$ with probability at least $2/3$
on every input if and only if it halts within time $T$ and outputs that indicator with
probability at least the constant function $2/3$ on every input.
\end{theorem}

\begin{definition}[Polynomial-time decision with success thresholds]
\lean{Randomized.BPP_PTMWith}
\statementsource{Vad12}{pdf-bf2ce5422088-p023-b001}
\statementsource{arora-barak-ch07-randomized-computation}{pdf-26d95080dbde-p002-b003}
\label{Randomized.BPP_PTMWith}
\leanok
\uses{Turing.FinNDTM, Turing.FinNDTM.BPDecidesInTimeWith}
For a function $p\colon\mathbb{N}\to\mathbb{Q}$, the class
$\mathrm{BPP}_{\mathrm{PTM}}(p)$ consists of languages $L$ over the Boolean alphabet for
which there are an exponent $e\in\mathbb{N}$, a constant $K\in\mathbb{N}$, and a
finite-state nondeterministic Turing machine that decides $L$ within time $K(n^e+1)$
with success threshold $p(n)$ on every input of length $n$.
\end{definition}

\begin{theorem}[Standard and threshold-based BPP coincide]
\lean{Randomized.BPP_PTM_eq_BPP_PTMWith}
\label{Randomized.BPP_PTM_eq_BPP_PTMWith}
\leanok
\uses{Randomized.BPP_PTM, Randomized.BPP_PTMWith}
\difficulty{1}
The class of languages decidable in polynomial time by a finite-state randomized Turing
machine with probability at least $2/3$ of producing the correct answer equals the class
defined by requiring that same lower bound on the probability of producing the correct
answer for every input.
\end{theorem}

\begin{definition}[Verifier characterization with success thresholds]
\lean{Randomized.InBPPWith}
\statementsource{Vad12}{pdf-bf2ce5422088-p023-b001}
\statementsource{arora-barak-ch07-randomized-computation}{pdf-26d95080dbde-p002-b003}
\label{Randomized.InBPPWith}
\leanok
\uses{Randomized.boolVerifier, Randomized.polyLen, Randomized.polyTimeModel, Randomized.randProb}
For a function $p\colon\mathbb{N}\to\mathbb{Q}$ and a Boolean language $L$, there are a
verifier $M$ and natural numbers $a,k$ such that the verifier has polynomial-time
efficiency, and for every Boolean input $x$, if $x\in L$ then the probability over
random strings $r$ of $M(x,r)=\mathrm{true}$ is at least $p(|x|)$, while if $x\notin L$
then the probability of $M(x,r)=\mathrm{false}$ is at least $p(|x|)$.
\end{definition}

\begin{theorem}[Fixed-error and standard BPP coincide]
\lean{Randomized.inBPPWith_iff_inBPP}
\label{Randomized.inBPPWith_iff_inBPP}
\leanok
\uses{Randomized.InBPP, Randomized.InBPPWith, Randomized.polyTimeModel}
\difficulty{1}
For every language $L$ of bit strings, $L$ has a polynomial-time verifier whose
probability of the correct answer on inputs of length $n$ is at least $2/3$ if and only
if it has a polynomial-time verifier whose probability of the correct answer is at least
$2/3$ on every input.
\end{theorem}

\begin{theorem}[Nondeterministic polynomial time implies BPP]
\lean{Randomized.inBPPWith_of_mem_BPP_PTMWith}
\label{Randomized.inBPPWith_of_mem_BPP_PTMWith}
\leanok
\uses{Complexity.exists_choiceVerifier_mem_P, Randomized.BPP_PTMWith, Randomized.InBPPWith, Randomized.append_mem_iff_outputWith_take, Randomized.eff_append_mem, Randomized.indicator_eq_false_iff, Randomized.polyLen, Randomized.randProb, Randomized.randProb_choiceVerifier, Randomized.randProb_mono, Randomized.randProb_take, Turing.FinNDTM.outputProb, Turing.FinNDTM.outputWith, Turing.MultiTapeTM.indicator, Turing.choiceBudget_le}
\difficulty{5}
For every rational-valued function $p$ on input lengths and every language $L$ of bit
strings, if $L$ has a nondeterministic polynomial-time machine that halts on every
choice string of the prescribed polynomial length and outputs the correct membership bit
with probability at least $p(n)$, then there is a polynomial-time randomized verifier
that outputs the correct membership bit with probability at least $p(n)$ on every input
of length $n$.
\end{theorem}

\begin{theorem}[From polynomial-time randomized machines to verifiers]
\lean{Randomized.inBPP_of_mem_BPP_PTM}
\label{Randomized.inBPP_of_mem_BPP_PTM}
\leanok
\uses{Randomized.BPP_PTM, Randomized.InBPP, Randomized.inBPPWith_of_mem_BPP_PTMWith, Randomized.polyTimeModel}
\difficulty{1}
Let $L$ be a language of bit strings that is decided by a finite nondeterministic Turing
machine with some running-time bound $c(n^d+1)$, where $c,d$ are natural numbers: for
every input of length $n$, every choice string of that length causes the machine to halt
within the bound, and the probability that its output is the correct membership bit is
at least $2/3$. Then there is a Boolean randomized verifier whose deterministic
acceptance and rejection relations are polynomial-time decidable and which, on each
input $x$, uses $a(|x|+1)^k$ fair random bits for some natural numbers $a,k$, outputting
the correct membership bit with probability at least $2/3$.
\end{theorem}

\begin{theorem}[Efficient machine for masked selection]
\lean{Randomized.exists_FinNDTM_maskSelect_of_eff}
\label{Randomized.exists_FinNDTM_maskSelect_of_eff}
\leanok
\uses{Complexity.Meyer.solveSplit_self', Complexity.P, Complexity.PolyTimeComputable, Complexity.exists_maskSelect_host, Complexity.maskSelect, Complexity.preimage_mem_P, Randomized.boolVerifier, Randomized.length_maskSelect, Randomized.polyLen, Randomized.polyTimeModel, Turing.FinNDTM, Turing.FinNDTM.outputWith, Turing.FinTM.computesFunInTime_splitSolve, Turing.MultiTapeTM.indicator, Turing.NDTM.initCfg, Turing.NDTM.runWith, Turing.pairEncode, Turing.solveSplit}
\difficulty{6}
Let $M$ map pairs of bit strings to a bit, and suppose the verifier that returns
$M(x,r)$ is efficient in the specified polynomial-time model. For any natural numbers
$a,k$, there are natural numbers $K,e$ and a finite nondeterministic Turing machine $N$
such that, for every input $x$, there are a natural number $s$ and a Boolean mask whose
number of true entries is $a(|x|+1)^k$, with $s+|\mathrm{mask}|\le K(|x|^e+1)$. On every
choice word $w$ of length exactly $K(|x|^e+1)$, $N$ halts, outputs the one-bit word
whose bit is $M(x,u)$ for $u$ obtained by applying the mask to the
length-$|\mathrm{mask}|$ segment of $w$ beginning at position $s$, and that selected
word has length $a(|x|+1)^k$.
\end{theorem}

\begin{theorem}[Efficient machine for randomized computation]
\lean{Randomized.exists_FinNDTM_of_eff}
\label{Randomized.exists_FinNDTM_of_eff}
\leanok
\uses{Complexity.maskSelect, Randomized.boolVerifier, Randomized.exists_FinNDTM_maskSelect_of_eff, Randomized.exists_ofFn_eq, Randomized.polyLen, Randomized.polyTimeModel, Randomized.randProb, Randomized.randProb_congr, Randomized.randProb_maskSelect_drop_take, Turing.FinNDTM, Turing.FinNDTM.outputProb, Turing.FinNDTM.outputWith, Turing.NDTM.HaltsWithin}
\difficulty{4}
Let $M$ map pairs of bit strings to a bit, and suppose the verifier that returns
$M(x,r)$ is efficient in the specified polynomial-time model. For any natural numbers
$a,k$, there are natural numbers $K,e$ and a finite nondeterministic Turing machine $N$
such that, for every input $x$, on every choice word of length exactly $K(|x|^e+1)$ the
machine halts; its output is a one-bit word; and for each bit $b$, the probability that
its output is $b$ equals the probability that $M(x,u)=b$ when $u$ is uniformly chosen
from bit strings of length $a(|x|+1)^k$. Moreover, on every such choice word, the output
equals $M(x,u)$ for some bit string $u$ of that length.
\end{theorem}

\begin{theorem}[Verifier characterization implies randomized polynomial time]
\lean{Randomized.mem_BPP_PTMWith_of_inBPPWith}
\label{Randomized.mem_BPP_PTMWith_of_inBPPWith}
\leanok
\uses{Randomized.BPP_PTMWith, Randomized.InBPPWith, Randomized.exists_FinNDTM_of_eff, Turing.MultiTapeTM.indicator}
\difficulty{4}
Let $p:\mathbb N\to\mathbb Q$ and let $L$ be a language of bit strings. If there are
constants $a,k\in\mathbb N$ and a polynomial-time Boolean verifier whose acceptance
probability on each $x\in L$ and rejection probability on each $x\notin L$, using
uniformly random bit strings of length $a(|x|+1)^k$, are each at least $p(|x|)$, then a
finite nondeterministic Turing machine decides $L$ with that error bound in polynomial
time, halting on every choice string of the prescribed length.
\end{theorem}

\begin{theorem}[Characterization of randomized polynomial time]
\lean{Randomized.mem_BPP_PTMWith_iff}
\label{Randomized.mem_BPP_PTMWith_iff}
\leanok
\uses{Randomized.BPP_PTMWith, Randomized.InBPPWith, Randomized.inBPPWith_of_mem_BPP_PTMWith, Randomized.mem_BPP_PTMWith_of_inBPPWith}
\difficulty{1}
For any function $p:\mathbb N\to\mathbb Q$ and any language $L$ of bit strings, the
following are equivalent. First, there are constants $K,e\in\mathbb N$ and a finite
nondeterministic Turing machine that, on each input $x$, halts on every choice string of
length $K(|x|^e+1)$ and outputs the membership indicator of $L$ with probability at
least $p(|x|)$. Second, there are constants $a,k\in\mathbb N$ and a polynomial-time
Boolean verifier such that, for every input $x$, its acceptance probability over
uniformly random bit strings of length $a(|x|+1)^k$ is at least $p(|x|)$ when $x\in L$,
and its rejection probability is at least $p(|x|)$ when $x\notin L$.
\end{theorem}

\begin{theorem}[Polynomial-Time Bounded-Error Recognition]
\lean{Randomized.mem_BPP_PTM_of_inBPP}
\label{Randomized.mem_BPP_PTM_of_inBPP}
\leanok
\uses{Randomized.BPP_PTM, Randomized.InBPP, Randomized.mem_BPP_PTMWith_of_inBPPWith, Randomized.polyTimeModel}
\difficulty{1}
Let $L$ be a language of bit strings. If there is a polynomial-time verifier that, on an
input string and a uniformly random bit string of polynomial length, accepts with
probability at least $2/3$ on inputs in $L$ and rejects with probability at least $2/3$
on inputs outside $L$, then $L$ belongs to $	extsf{BPP}_{
m PTM}$.
\end{theorem}

\begin{theorem}[Characterization of bounded-error polynomial time]
\lean{Randomized.mem_BPP_PTM_iff}
\label{Randomized.mem_BPP_PTM_iff}
\leanok
\uses{Randomized.BPP_PTM, Randomized.InBPP, Randomized.mem_BPP_PTMWith_iff, Randomized.polyTimeModel}
\difficulty{1}
For every language $L$ of bit strings, membership in the class of languages decided by
finite nondeterministic Turing machines that halt on every random choice string of
polynomial length and output the correct membership indicator with probability at least
$2/3$ is equivalent to the existence of a polynomial-time Boolean verifier with
polynomially many uniformly random bits that accepts members and rejects nonmembers with
probability at least $2/3$.
\end{theorem}

\begin{theorem}[Machine and verifier definitions of BPP agree]
\lean{Randomized.BPP_PTM_eq}
\label{Randomized.BPP_PTM_eq}
\leanok
\uses{Randomized.BPP_PTM, Randomized.InBPP, Randomized.mem_BPP_PTM_iff, Randomized.polyTimeModel}
\difficulty{1}
For the specified polynomial-time verifier model, the class of languages decidable in
polynomial time by a finite-state randomized Turing machine with probability at least
$2/3$ of producing the correct answer equals the class for which an efficient Boolean
verifier, given a polynomial number of uniformly random bits, outputs the correct answer
with probability at least $2/3$ on every input.
\end{theorem}

\begin{definition}[One-sided randomized time-bounded decision]
\lean{Turing.FinNDTM.RDecidesInTimeWith}
\statementsource{Vad12}{pdf-bf2ce5422088-p021-b006}
\statementsource{arora-barak-ch07-randomized-computation}{pdf-26d95080dbde-p006-b003}
\label{Turing.FinNDTM.RDecidesInTimeWith}
\leanok
\uses{Turing.FinNDTM, Turing.FinNDTM.outputProb, Turing.FinNDTM.outputWith, Turing.NDTM.HaltsWithin}
Let $N$ be a finite-state nondeterministic Turing machine over the Boolean alphabet, let
$L$ be a language of Boolean strings, let $T:bn	obn$ be a time bound, and let
$p:bn	obq$. For every input string $x$, the machine halts within $T(|x|)$ steps; if $x
otin L$, every choice string of length exactly $T(|x|)$ makes its output the one-symbol
string consisting of false; and if $x\in L$, the probability that its output is the
one-symbol string consisting of true after $T(|x|)$ steps is at least $p(|x|)$.
\end{definition}

\begin{theorem}[Equivalent randomized time criteria]
\lean{Turing.FinNDTM.rDecidesInTime_iff}
\label{Turing.FinNDTM.rDecidesInTime_iff}
\leanok
\uses{Turing.FinNDTM, Turing.FinNDTM.RDecidesInTime, Turing.FinNDTM.RDecidesInTimeWith}
\difficulty{1}
Let $N$ be a finite nondeterministic Turing machine over the binary alphabet, let $L$ be
a language of binary strings, and let $T:\mathbb{N}\to\mathbb{N}$. The following are
equivalent: for every binary string $x$, every run of length $T(|x|)$ halts, members
$x\in L$ have probability at least $2/3$ of producing the one-bit output $1$, and
nonmembers produce the one-bit output $0$ on every choice word of length $T(|x|)$; and
the same conditions with the member probability bound written as $2/3\leq
\Pr[\text{output }1]$.
\end{theorem}

\begin{definition}[Polynomial-time randomized machine class]
\lean{Randomized.RP_PTMWith}
\statementsource{Vad12}{pdf-bf2ce5422088-p021-b006}
\statementsource{arora-barak-ch07-randomized-computation}{pdf-26d95080dbde-p006-b003}
\label{Randomized.RP_PTMWith}
\leanok
\uses{Turing.FinNDTM, Turing.FinNDTM.RDecidesInTimeWith}
For a function $p:bn\to\bbq$, let $\mathrm{RP}_{p}$ be the union, over all exponents
$e\in\bbn$, of the languages $L$ for which there are a constant $K\in\bbn$ and a
finite-state nondeterministic Turing machine over Boolean symbols that decides $L$ with
the preceding one-sided error guarantees within time $K(n^e+1)$ on inputs of length $n$,
using threshold $p(n)$.
\end{definition}

\begin{theorem}[Equivalent Definitions of One-Sided Randomized Polynomial Time]
\lean{Randomized.RP_PTM_eq_RP_PTMWith}
\label{Randomized.RP_PTM_eq_RP_PTMWith}
\leanok
\uses{Randomized.RP_PTM, Randomized.RP_PTMWith}
\difficulty{1}
The class of languages decidable in polynomial time by a randomized machine with
one-sided error and acceptance probability at least $2/3$ on members equals the class
defined in the same way using the fixed member-input acceptance threshold $2/3$.
\end{theorem}

\begin{definition}[Randomized verifier characterization]
\lean{Randomized.InRPWith}
\statementsource{Vad12}{pdf-bf2ce5422088-p021-b006}
\statementsource{arora-barak-ch07-randomized-computation}{pdf-26d95080dbde-p006-b003}
\label{Randomized.InRPWith}
\leanok
\uses{Randomized.boolVerifier, Randomized.polyLen, Randomized.polyTimeModel, Randomized.randProb}
Let $p:\bbn\to\bbq$ and let $L$ be a language of Boolean strings. There exist a
Boolean-valued verifier $M(x,r)$ and natural numbers $a,k$ such that the verifier runs
in polynomial time, every $x\in L$ is accepted with probability at least $p(|x|)$ when
given a uniformly random Boolean string of length $\mathrm{polyLen}(a,k,|x|)$, and every
$x\notin L$ is rejected for every Boolean string of that length.
\end{definition}

\begin{theorem}[Characterization of One-Sided Randomized Time]
\lean{Randomized.mem_RP_PTMWith_iff}
\label{Randomized.mem_RP_PTMWith_iff}
\leanok
\uses{Complexity.exists_choiceVerifier_mem_P, Randomized.InRPWith, Randomized.RP_PTMWith, Randomized.append_mem_iff_outputWith_take, Randomized.eff_append_mem, Randomized.exists_FinNDTM_of_eff, Randomized.indicator_eq_false_iff, Randomized.polyLen, Randomized.randProb_choiceVerifier, Turing.MultiTapeTM.indicator, Turing.choiceBudget_le}
\difficulty{5}
For every function $p:\mathbb{N}\to\mathbb{Q}$ and every language $L$ of bit strings,
$L$ belongs to the class of languages decided in polynomial time with acceptance
probability at least $p(|x|)$ on members and certain rejection on nonmembers exactly
when there is a polynomial-time verifier using polynomially many random bits that has
those same acceptance and rejection guarantees.
\end{theorem}

\begin{theorem}[Characterization of Polynomial-Time RP]
\lean{Randomized.mem_RP_PTM_iff}
\label{Randomized.mem_RP_PTM_iff}
\leanok
\uses{Randomized.InRP, Randomized.RP_PTM, Randomized.RP_PTM_eq_RP_PTMWith, Randomized.mem_RP_PTMWith_iff, Randomized.polyTimeModel}
\difficulty{3}
For every language $L$ of bit strings, $L$ belongs to the polynomial-time RP class
exactly when there is a polynomial-time verifier using polynomially many uniformly
random bits that accepts members with probability at least $2/3$ and rejects nonmembers
on every random string.
\end{theorem}

\begin{theorem}[Characterization of coRP with Polynomial-Time Verifiers]
\lean{Randomized.mem_coRP_PTM_iff}
\label{Randomized.mem_coRP_PTM_iff}
\leanok
\uses{Randomized.InCoRP, Randomized.coRP_PTM, Randomized.mem_RP_PTM_iff, Randomized.polyTimeModel}
\difficulty{1}
A language $L$ belongs to coRP with polynomial-time machines if and only if its
complement has a polynomial-time verifier with one-sided error: for some natural numbers
$a,k$, it uses $a(|x|+1)^k$ random bits, accepts each input in the complement with
probability at least $2/3$, and rejects every input outside the complement on every
random string. The verifier is polynomial-time efficient in the sense that its
true-output and false-output sets on encoded input-randomness pairs both lie in $P$.
\end{theorem}

\begin{theorem}[Bpp ptm subset exp]
\lean{Randomized.BPP_PTM_subset_EXP}
\statementsource{arora-barak-ch07-randomized-computation}{pdf-26d95080dbde-p003-b002}
\label{Randomized.BPP_PTM_subset_EXP}
\leanok
\uses{Complexity.EXP, Randomized.BPP_PTM, Randomized.inBPP_of_mem_BPP_PTM, Randomized.mem_EXP_of_inBPP_polyTimeModel}
\difficulty{1}
% dataset-pipeline: stub -- prose pending build_informal_statements.py
Statement pending regeneration.
\end{theorem}

\begin{theorem}[BPP equals P when P equals NP]
\lean{Randomized.BPP_PTM_eq_P_of_P_eq_NP}
\label{Randomized.BPP_PTM_eq_P_of_P_eq_NP}
\leanok
\uses{Complexity.NP, Complexity.P, Randomized.BPP_PTM, Randomized.inBPP_polyTimeModel_iff_mem_P_of_P_eq_NP, Randomized.mem_BPP_PTM_iff}
\difficulty{3}
If $\mathrm{P}=\mathrm{NP}$, then $\mathrm{BPP}=\mathrm{P}$.
\end{theorem}

\begin{theorem}[Polynomial randomized time and polynomial circuit size]
\lean{Randomized.BPP_PTM_subset_PPoly}
\label{Randomized.BPP_PTM_subset_PPoly}
\leanok
\uses{BoolCircuit.PPoly, BoolCircuit.mem_PPoly_iff, Randomized.BPP_PTM, Randomized.adleman_polyTime, Randomized.inBPP_of_mem_BPP_PTM}
\difficulty{2}
Every language decidable by a bounded-error randomized Turing machine in polynomial time
is decided by a family of Boolean circuits of fan-in at most two whose sizes are bounded
by a polynomial in the input length.
\end{theorem}

\begin{theorem}[Randomized polynomial time and the second hierarchy level]
\lean{Randomized.BPP_PTM_subset_SigmaP_two_inter_PiP_two}
\label{Randomized.BPP_PTM_subset_SigmaP_two_inter_PiP_two}
\leanok
\uses{Complexity.PiP, Complexity.SigmaP, Randomized.BPP_PTM, Randomized.inBPP_of_mem_BPP_PTM, Randomized.sipser_gacs_polyTime}
\difficulty{1}
Every language decidable by a bounded-error randomized Turing machine in polynomial time
belongs both to the second existential level of the polynomial hierarchy and to the
second universal level; equivalently, it belongs to their intersection.
\end{theorem}

\begin{theorem}[BPP is closed under polynomial-time reductions]
\lean{Randomized.BPP_PTM_of_polyTimeReducible}
\label{Randomized.BPP_PTM_of_polyTimeReducible}
\leanok
\uses{Randomized.BPP_PTM, Randomized.InBPP.of_polyTimeReducible, Randomized.inBPP_of_mem_BPP_PTM, Randomized.mem_BPP_PTM_of_inBPP}
\difficulty{1}
Let $L,L'$ be languages of bit strings. If $L'$ is polynomial-time many-one reducible to
$L$ and $L$ belongs to BPP, then $L'$ belongs to BPP.
\end{theorem}

\begin{theorem}[A Collapse of the Polynomial Hierarchy]
\lean{Randomized.PH_eq_SigmaP_two_of_SAT3_mem_BPP_PTM}
\label{Randomized.PH_eq_SigmaP_two_of_SAT3_mem_BPP_PTM}
\leanok
\uses{Complexity.PH, Complexity.SAT3, Complexity.SigmaP, Randomized.BPP_PTM, Randomized.PH_eq_SigmaP_two_of_inBPP_SAT3, Randomized.inBPP_of_mem_BPP_PTM}
\difficulty{1}
If $3$-CNF satisfiability belongs to the polynomial-time bounded-error probabilistic
class $\mathrm{BPP}_{\mathrm{PTM}}$, then the polynomial hierarchy equals its second
existential level: $\mathrm{PH}=\Sigma_2^P$.
\end{theorem}

\begin{theorem}[Error Reduction for One-Sided Randomized Polynomial Time]
\lean{Randomized.RP_PTMWith_error_reduction}
\label{Randomized.RP_PTMWith_error_reduction}
\leanok
\uses{Randomized.RP_PTMWith, Randomized.mem_RP_PTMWith_iff, Randomized.rp_error_reduction_polyTime}
\difficulty{1}
For all natural numbers $c$ and $d$, every language decidable in polynomial time by a
randomized machine with one-sided error and acceptance probability on members at least
$\min\!\left(\frac{2}{3},(n+1)^{-c}\right)$ on inputs of length $n$ is decidable in
polynomial time by a randomized machine with one-sided error and acceptance probability
on members at least $1-2^{-((n+1)^d+1)}$.
\end{theorem}
